\section{Introduction}
\label{sec:intro}

A time-series foundation model is pretrained once on a large heterogeneous corpus, then asked to forecast an
unseen series with no fitting, no feature engineering and no local expertise. If that promise holds where
forecasting capacity is scarcest, it matters. Those are the places with the fewest stations, the fewest trained
forecasters and the least compute, and they are often where a week of warning carries the largest consequences.
A model that runs zero-shot on a free GPU is precisely the technology they lack.

Whether it holds there is an empirical question current benchmarks do not answer, because the evidence base is
skewed in two ways at once. Geographically, general time-series suites treat weather as one domain among many and
draw it overwhelmingly from mid-latitude sources \citep{godahewa2021monash,aksu2024gifteval,qiu2024tfb,li2025tsfmbench};
across seven widely used benchmark papers, the words ``monsoon'' and ``Bangladesh'' do not appear. Representationally,
the weather benchmarks that \emph{are} globally comprehensive score against reanalysis
\citep{rasp2024weatherbench2,nathaniel2024chaosbench}. Reanalysis is a physically consistent
reconstruction, smoother than the station record it interpolates; it is smoothest exactly where stations are
sparsest. Observation-based
evaluation exists for numerical and AI weather-prediction models \citep{jin2024weatherreal,gupta2026mausam}, but the
corresponding question for zero-shot time-series models has not been asked.

Bangladesh makes the question concrete. Its weather is governed by the South Asian summer monsoon, whose defining
feature, the alternation of active and break spells inside the rainy season, has no temperate analogue. The 35-station
record we build on shows the range. Median July rainfall reaches \SI{457}{\milli\metre} per station-month, against
\SI{0}{\milli\metre} in January. \SI{67.6}{\percent} of station-days are dry, yet the wet-day 99th percentile
is \SI{143}{\milli\metre}. The 2016--2023 test period alone contains 432 cyclone station-days across 14 storms. The
record that would support decisions here is thin. Over the same period the ten Bangladeshi stations in GHCN-Daily
are about \SI{65}{\percent} complete, with 2021 almost entirely absent. The European stations we use as a
reference are \SI{99}{\percent} complete. That asymmetry is not only a property of some pretraining corpus: it sits inside
the model's input window at forecast time.

We introduce \textbf{BanglaWeatherBench}, a benchmark designed so that its headline number is not the only thing it
produces. It defines 16 daily forecasting tasks over 170{,}571 rolling-origin windows across four tracks. The
forecast target is 35 Bangladesh Meteorological Department stations; beside them sit 10 co-located GHCN-Daily
stations, 32 European GHCN-Daily stations as a temperate reference, and NASA POWER reanalysis at all 77 station
coordinates. A seventeenth task covers CHIRPS dekadal rainfall over 64 districts, from an independent source. Three controls answer the confounds that would otherwise make a tropical
result unreadable. The same-network temperate reference isolates climatic regime; the reanalysis track
at identical coordinates isolates the effect of the scoring target. The gap-matched context control isolates data
availability from climate, and we apply it to the trained baselines as well as to the foundation models. Scoring is
stratified by monsoon phase and by active/break spell.

Two commitments shape what we are willing to claim. First, significance. A 7-day stride with a 30-day horizon
produces heavily overlapping windows, so we resample at the forecast origin. We report Diebold--Mariano tests with
HAC variances, moving-block bootstrap intervals, Friedman--Nemenyi ranks, and Holm correction across 48 comparisons.
Many margins that look decisive as point estimates do not survive. Second, leakage: we audit each model's
pretraining corpus against its technical report. Leakage-free zero-shot evaluation is not itself novel
\citep{qiao2026time}; here it is a precondition for reading a geographic contrast as a geographic effect.

Evaluating 8 zero-shot foundation models against 10 trained and statistical baselines yields a mixed result that we
think is more useful than either a triumphal or a dismissive one.

\begin{enumerate}
  \item \textbf{Foundation models lead, but most of the lead is not separable.} The best model beats the best
    baseline in 14 of 48 Holm-corrected comparisons, loses 5 and is indistinguishable in 29; on station
    observations it wins 3 of 24, all at a one-day lead, against 11 of 24 on reanalysis. The six leading models and
    PatchTST sit within one critical difference.
  \item \textbf{Nothing beats climatology on tropical rainfall} beyond a one-day lead at the BMD stations, or at
    any horizon on dekadal rainfall, a concrete boundary on the ``forecast anything'' claim.
  \item \textbf{Reanalysis flatters models} relative to the stations it represents, in 78 of 108 comparisons.
  \item \textbf{The tropical temperature deficit is narrow and mostly about data.} At a 30-day lead 4 of 8 models
    lose significantly more skill than a like-for-like control; about \SI{60}{\percent} of the raw penalty is
    explained by sparser context, and doubling the context does not remove the remainder.
  \item \textbf{Skill reverses by regime}: strong in break spells, significantly worse than climatology in active
    spells, in 61 of 72 tests.
\end{enumerate}

We state finding 4 carefully because it is easy to misreport. The evidence supports a specific claim: for
temperature at long horizons, half the roster loses significantly more skill on Bangladeshi stations than trained
baselines do, most of it traceable to sparser context. It does not support a blanket claim that foundation
models fail in the Global South. Stated correctly it is also the more actionable result: much of the deficit is a
data-availability problem that observation programmes and gap-tolerant architectures can address.

\paragraph{Contributions.} (i) A released, leakage-tested benchmark for a tropical-monsoon climate, with splits,
strata, a datasheet and per-source licensing. (ii) Controls rather than a leaderboard alone: a same-network
temperate reference, a reanalysis track at identical coordinates, and a gap-matched control applied to models and
baselines alike. (iii) A significance protocol matched to overlapping rolling-origin panels, applied uniformly.
(iv) To our knowledge the first evaluation of current zero-shot time-series foundation models on South Asian
station observations including precipitation, stratified by monsoon phase, with a pretraining-overlap audit.
(v) Reproducibility at the compute scale of the places the benchmark is about: 4.16 GPU-hours on a free Tesla T4.
